\documentclass[11pt,a4paper]{article}
\usepackage[margin=26mm]{geometry}
\usepackage{amsmath,amssymb,amsthm,microtype,hyperref}
\newcommand{\E}{\mathbb E}\newcommand{\Pp}{\mathbb P}
\newtheorem{theorem}{Theorem}\newtheorem{lemma}[theorem]{Lemma}
\title{Structural count bounds beyond the small-noise comparison}
\author{Research note, 28 October 2026}\date{}
\begin{document}\maketitle

The direct comparison through $ns=.075$ is retained unchanged. This note
supplies two tools for the remaining noise range: a concentration bound
from the ordered conditional counts, and an analytic upper bound for the
signed binomial curvature correction. The first gives a complete comparison
in the symmetric central subcase with no noise ceiling. The tools do not
yet close the general larger-noise comparison.

\section{Concentration from the ordered counts}
Let $Y=B+U$, where $\Pp(B=1)=p\in(0,1)$, $\E(U\mid B)=0$, and
\[
 U_0\le a_0,\qquad U_1\ge-a_1,\qquad
 a_0,a_1\ge0,\quad \zeta=1-a_0-a_1>0.
\]
Take independent copies $U_{0i},U_{1i}$ and set
\[
 T_j=j+\sum_{i\le j}U_{1i}+\sum_{i>j}U_{0i},\qquad
 \pi_j=\binom njp^j(1-p)^{n-j},\quad0\le j\le n.
\]
Every conditional count has the correct law, and
$T_{j+1}-T_j\ge\zeta$ on each coupled sample.

\begin{theorem}\label{thm:concentration}
For a closed interval $I$ of length less than $m\zeta$, where
$1\le m\le n+1$ is an integer,
\[
 \Pp(Y_1+\cdots+Y_n\in I)
       \le\max_{0\le j\le n+1-m}\sum_{r=j}^{j+m-1}\pi_r.
\]
In particular every atom, and every interval of length less than $\zeta$,
has probability at most $\max_j\pi_j$.
\end{theorem}
\begin{proof}
For each coupled sample, the indices with $T_j\in I$ are consecutive,
and at most $m$ of them can occur: $m+1$ points would have span at least
$m\zeta$. Conditioning on the actual binomial count gives
\[
 \Pp(Y_1+\cdots+Y_n\in I)
       =\E\sum_{j=0}^n\pi_j\mathbf1_{\{T_j\in I\}}.
\]
A shorter block can be extended to $m$ legal indices, so its mass is
bounded by the displayed maximum. This proves the claim without any
independence assertion about a selected index. The strict length condition
matters: for noiseless Bernoulli variables, $[j,j+1]$ contains two atoms.
\end{proof}

\begin{theorem}[Symmetric central comparison at every noise level]
Let $n$ be even, $p=1/2$, $U_1$ have the law of $-U_0$,
$\E U_0=0$, $U_0\le a<1/2$, and assume finite third absolute moments.
Set $s=\E U_0^2>0$, $X=(Y-1/2)/\sqrt{1/4+s}$ and
$\rho=\E|X|^3$. For every $C>0$ and
\[
             0\le H<\frac{1-2a}{2\sqrt{1/4+s}},
\]
the direct comparison with the central Bernoulli anchor $k=n/2$ is strict:
\begin{align*}
 &\sqrt n\{\Pp(X_1+\cdots+X_n\le H)-\Phi(H/\sqrt n)\}-C\rho\\
 &\hspace{15mm}-\sqrt n\{F_B(n/2)-1/2\}+C<0.
\end{align*}
There is no upper bound on $ns$ in this statement.
\end{theorem}
\begin{proof}
The marginal $X$ is symmetric. For $H\ge0$ this gives the exact identity
\[
 \Pp(X_1+\cdots+X_n\le H)
 =\frac12+\frac12\Pp(|X_1+\cdots+X_n|\le H).
\]
The corresponding raw interval for the sum of the $Y_i$ has length
$2H\sqrt{1/4+s}<1-2a$. Theorem~\ref{thm:concentration} bounds its
mass by the central binomial mass $\pi_{n/2}$. Also
$F_B(n/2)=1/2+\pi_{n/2}/2$. Subtraction therefore bounds the displayed
comparison by
\[
       -\sqrt n\{\Phi(H/\sqrt n)-1/2\}-C(\rho-1).
\]
Jensen's strict inequality gives $\rho>1$ unless $|X|=1$ almost surely.
The two conditional supports lie strictly on opposite sides of zero,
since $a<1/2$. Thus $|X|=1$ would force each conditional law to be a
point mass and $s=0$, contrary to the hypothesis. The final bound is strict,
including at $H=0$ and at exact atom collisions.
\end{proof}

\section{The signed curvature correction has a short uniform bound}
Write $v_j=\E U_j^2$, $q=1-p$, $v=pq$, $s=qv_0+pv_1$, and
$W_j=T_j-j$. For any real $t$, define
\[
 B_j(t)=\E(t-T_j)_+-(t-j)_+\ge0,\qquad
 \mathcal C(t)=\sum_{j=1}^{n-1}(\pi_{j+1}-\pi_j)B_j(t).
\]
These quantities require only finite second moments. The new crossing
identity in the companion \texttt{crossing.tex} writes the actual CDF as
the linear interpolation of the binomial CDF, plus $\mathcal C(t)$,
minus a weighted nonnegative crossing loss. The following estimate bounds
only the correction; it is not a bound for that loss.

\begin{lemma}\label{lem:curvature}
For $n\ge2$, $0<p<1$ and every real $t$, put $c=(np-q-t)_+$.
Then
\begin{equation}\label{eq:curvature}
 \mathcal C(t)\le\frac{s}{4v}
             +\frac{c\sqrt s}{2v\sqrt{n+1}}.
\end{equation}
In terms of $\lambda=ns$, this is
\[
 \sqrt n\,\mathcal C(t)\le\frac{\lambda}{4v\sqrt n}
                +\frac{c\sqrt\lambda}{2v\sqrt{n+1}}.
\]
The second term vanishes whenever $t\ge np-q$.
\end{lemma}
\begin{proof}
Let $a=t-j$ and $t_j=(n-j)v_0+jv_1$. Centering gives
\[
 B_j(t)=\frac{\E|a-W_j|-|a|}{2}.
\]
The triangle inequality and Cauchy--Schwarz imply
$B_j\le\sqrt{t_j}/2$. When $a>0$, another application of
Cauchy--Schwarz gives
\[
 B_j\le\frac{\sqrt{a^2+t_j}-a}{2}\le\frac{t_j}{4a}.
\]
Thus $a_+B_j\le t_j/4$ for all $a$.
Write $\widetilde\pi_r=\Pp(\operatorname{Bin}(n+1,p)=r)$.
The exact adjacent-mass identity is
\[
 \pi_{j+1}-\pi_j
   =\frac{(n+1)p-j-1}{(n+1)v}\widetilde\pi_{j+1}.
\]
Its numerator is $a+(np-q-t)\le a_++c$. Discard the nonpositive
mass differences, using $B_j\ge0$. On the positive side,
\[
 \mathcal C(t)\le\frac1{(n+1)v}
   \sum_j\widetilde\pi_{j+1}
        \left(\frac{t_j}{4}+\frac c2\sqrt{t_j}\right).
\]
The remaining sums may be extended over all $1\le j\le n-1$.
For $r=j+1$, $t_j=(n+1-r)v_0+(r-1)v_1$ is at most
$(n+1-r)v_0+rv_1$. Extending this latter nonnegative expression over
the complete binomial distribution gives
\[
 \sum_{j=1}^{n-1}\widetilde\pi_{j+1}t_j\le(n+1)s.
\]
Cauchy--Schwarz also bounds the corresponding square-root sum by
$\sqrt{(n+1)s}$, since its total mass is at most one.
Substitution proves \eqref{eq:curvature}. No local truncation of the
binomial distribution, or numerical count grid, was used.
\end{proof}

The signed sum can be much smaller than this upper bound. The proof
retains its favorable sign long enough to discard only the decreasing
side of the mass function. Nevertheless, below $np-q$ the remaining
square-root term is not a bound proportional to arbitrarily small noise.
Even above $np-q$, a quantitative weighted crossing-loss estimate and the
normal/cubic repair are still needed to prove the full comparison.

\section{A precise obstruction to one four-moment shortcut}
The nonnegative local crossing kernel gives the sufficient lower bound
\[
 g_n(t)\ge\Pp(W_i>t-i)-\E(W_i-(t-i))_+
\]
by selecting one index. It would be useful if the expectation on the right
were positive from the usual own-count moment bounds alone. This is false.
Let
\[
 \Pp(Z=6/5)=\Pp(Z=-6/5)=25/144,
                 \qquad\Pp(Z=0)=47/72.
\]
Then
\[
 \E Z=\E Z^3=0,\quad \E Z^2=1/2,\quad\E Z^4=18/25
       \le(3/5)^2(1/2)+3(1/2)^2=93/100,
\]
but
\[
                   \Pp(Z>0)-\E Z_+=-5/144<0.
\]
Thus no polynomial argument using only these four moment constraints can
prove positivity of that selected reserve on the whole moderate-noise
range. This law is a counterexample to the moment relaxation, not an
asserted convolution of admissible iid conditional variables. Additional
independence information, joint neighboring terms, or the full nonnegative
kernel is necessary.

A second false shortcut concerns renewal independence. Take $A=3/5$,
$a=3/125$ and the comonotone conditional coupling
\[
 (U_0,U_1)=(-A,-a),(a,-a),(a,A)
       \quad\text{with probabilities }1/26,12/13,1/26.
\]
Both marginals are centered and obey the right strip with inward sum
$2a=.048$. At $n=42,p=1/2$ their noise is
$ns=42Aa=378/625=.6048$. Although the increment
$D=1+U_1-U_0$ has mean one and width $78/125<1$,
\[
                      \operatorname{Cov}(U_0,D)=-216/15625\ne0.
\]
The initial mesh location $\sum U_{0i}$ therefore remains correlated
with its increments. An ordinary independent-delay renewal theorem cannot
be applied merely by choosing this coupling. The companion exact
stop-loss telescoping identity accounts for the bias and avoids the issue.

\section{Evidence and the remaining inequality}
The exact checker uses actual centered conditional two-point marginals,
including one vanishing conditional noise and exact collisions. It verifies
interval concentration, the binomial difference identity, the weighted
variance step, the curvature bound and the symmetric central comparison.
Radical bounds in the curvature control are checked by squaring only after
checking the sign of the residual. The two counterexamples above are exact.

A separate bounded diagnostic tested 115 actual four/three-point laws at
$n=42,p=1/2$ and noise $.075<ns\le.616$. It evaluated the unweighted
crossing at 43,487 exact atom locations near the mean, with floating
probabilities; all observed losses were positive. Those observations are
discovery only and are not inputs to any theorem here.

The unresolved general implication is quantitative. In the exact identity
\[
 \Delta_C=-\sqrt n\sum_i\pi_i g_i(t)
       +\sqrt n\,\mathcal C(t)+\text{phase, normal and cubic repairs},
\]
one must bound the first, weighted loss strongly enough for the same
actual law to absorb every remaining term. The fact that
$\sum_i g_i(t)\ge0$, even for the whole acquired noise range, does not
supply such a bound. The existing accepted count certificates remain in
place; no larger-noise full theorem or global propagation is claimed.
\end{document}
